\section{Antecedentes y objetivos de la charla}
\subsection{El ponente y su lugar en el curso}
Esta charla pertenece al curso avanzado UC Berkeley CS294-280 Sp25 y la dirige Jason Weston, de Meta Research. Desde el inicio, la charla señala que la investigación actual en reasoning atraviesa saltos comparables a o1/R1 y DeepSeek-R1: los modelos avanzan con rapidez en los benchmark multimodales, pero cada salto también revela que los errores de razonamiento no se detectan a tiempo. Por eso, esta sección busca reconstruir el lugar de la charla: es a la vez un resumen de la historia reciente de Chain-of-Thought (que en realidad se popularizó hasta 2022) y un puente entre «dejar que el modelo defina por sí mismo la tarea de reasoning» y «dejar que él mismo juzgue el resultado».

\subsection{Propósito de la charla}
Weston resume el tema como self-improving language models: modelos que, durante el entrenamiento, aprenden poco a poco a formular sus propias preguntas, calificarse a sí mismos y optimizarse. Al inicio, el ponente menciona varias veces que «me parece que esta generación de investigadores de IA tiene mucha suerte de poder presenciar estos avances tan veloces» y subraya que «de forma análoga a un evento de benchmark como DeepSeek-R1, incluso en unos pocos meses se puede observar un cambio cualitativo en la reasoning accuracy». Por eso, a continuación las notas se desarrollan en torno a cuatro dimensiones: el objetivo del reasoning, el ciclo cerrado de entrenamiento, el sistema de evaluación y los retos del despliegue.

\begin{knowledgebox}{El contexto histórico de Chain-of-Thought}
El marco de entrenamiento de Chain-of-Thought se popularizó de manera explosiva a partir de 2022; Weston lo considera «un punto de partida rápido para el reasoning», pero también advierte que la verdadera capacidad de razonamiento exige que el modelo no solo produzca la cadena de razonamiento, sino que también sea capaz de evaluar cuál de las cadenas es más confiable.
\end{knowledgebox}

\subsection{Resumen de la sección}
Esta sección, a partir de hitos históricos (el auge de Chain-of-Thought, DeepSeek-R1, o1/R1) y de las reflexiones del ponente sobre el «self-improvement», establece el objetivo general y el criterio de evaluación de toda la charla. La pregunta central que hay que atender a continuación es: ¿cómo descomponer la «capacidad de introspección del modelo» en pasos concretos de entrenamiento?
